\section{Number of Operations to Make Network Connected}
\label{sec:graph-network-connected}

\problemstatement{Given $n$ computers and a list of existing cable
connections, find the \textbf{minimum number of operations} to make
every computer connected, where one operation removes a cable and uses
it to connect any two currently-disconnected computers. Return $-1$ if
there are not enough cables to possibly connect everything.}

\begin{patternbox}
\emph{``Reuse redundant connections to join separate groups''} is a
direct Disjoint Set application: union every given connection, then
count how many separate \textbf{components} remain. Connecting $k$
components into $1$ always takes exactly $k{-}1$ operations (each
operation merges two components into one), and a \textbf{redundant}
cable (one connecting two computers \emph{already} in the same
component) is exactly a ``spare'' cable available to be moved.
\end{patternbox}

\subsection*{Why redundant cables exactly cover the requirement}

A graph with $n$ nodes and $c$ connected components needs exactly
$c{-}1$ more edges to become fully connected (this is the same
spanning-tree counting argument from Section~\ref{sec:graph-kruskal}:
each component is already internally connected, like a single
``super-node,'' and connecting $c$ super-nodes into one tree needs
$c{-}1$ edges). The number of redundant cables (edges that connected
two nodes already in the same component when processed) is exactly
$\textit{totalCables} - (n - c)$, since a spanning structure over $n$
nodes split into $c$ components uses exactly $n-c$ non-redundant
edges. So the operation is feasible exactly when the redundant count
is $\ge c{-}1$ -- equivalently, exactly when $\textit{totalCables}
\ge n-1$.

\subsection*{Algorithm}

\begin{enumerate}
  \item If the number of connections is less than $n{-}1$, return
        $-1$ immediately (not enough cable to possibly connect $n$
        computers).
  \item Union every connection using Disjoint Set.
  \item Count the number of distinct components $c$ remaining (count
        nodes that are their own root).
  \item Return $c{-}1$.
\end{enumerate}

\subsection*{C++ implementation}

\begin{lstlisting}
#include <bits/stdc++.h>
using namespace std;

class DisjointSet {
    vector<int> parent, size;
public:
    DisjointSet(int n) : parent(n), size(n, 1) {
        for (int i = 0; i < n; i++) parent[i] = i;
    }
    int find(int x) {
        if (parent[x] == x) return x;
        return parent[x] = find(parent[x]);
    }
    void unite(int x, int y) {
        int rootX = find(x), rootY = find(y);
        if (rootX == rootY) return;
        if (size[rootX] < size[rootY]) swap(rootX, rootY);
        parent[rootY] = rootX;
        size[rootX] += size[rootY];
    }
};

int makeConnected(int n, vector<vector<int>>& connections) {
    if ((int)connections.size() < n - 1) return -1;

    DisjointSet ds(n);
    for (auto& c : connections) ds.unite(c[0], c[1]);

    int components = 0;
    for (int i = 0; i < n; i++) {
        if (ds.find(i) == i) components++;
    }
    return components - 1;
}

int main() {
    vector<vector<int>> connections = {{0,1},{0,2},{1,2}};
    cout << makeConnected(4, connections) << endl; // 1
    return 0;
}
\end{lstlisting}

\subsection*{Dry run}

\dryrun{Take $n=4$, $\textit{connections}=[[0,1],[0,2],[1,2]]$.}

$3$ connections $\ge n{-}1=3$: feasible. Union $0{-}1$, $0{-}2$
(now $\{0,1,2\}$ connected), $1{-}2$ (redundant -- already same
component). After processing, components: $\{0,1,2\}$ (root, say,
$0$) and $\{3\}$ (its own root) -- $c=2$. Answer: $c-1=1$ -- one
redundant cable (the $1{-}2$ connection) can be moved to connect node
$3$.

\begin{complexitybox}
\textbf{Time Complexity.} $O((n+E)\cdot\alpha(n))$ -- unioning every
connection plus a final scan over all nodes.
\textbf{Space Complexity.} $O(n)$ for the Disjoint Set.
\end{complexitybox}

\begin{takeawaybox}
This is the first of several problems in this chapter where Disjoint
Set is used purely for its \textbf{grouping} ability -- counting and
merging components -- with no notion of edge weight or minimality
involved at all, distinguishing it clearly from the MST algorithms
earlier in this chapter that also relied on Disjoint Set.
\end{takeawaybox}
